\documentclass[12pt,a4paper]{article}
\usepackage[utf8]{inputenc}
\usepackage[spanish]{babel}
\usepackage{geometry}
\geometry{margin=2.5cm}
\usepackage{helvet}
\renewcommand{\familydefault}{\sfdefault}
\usepackage{graphicx}
\usepackage{amsmath}
\usepackage{booktabs}
\usepackage{float}
\usepackage{caption}
\usepackage{enumitem}
\usepackage{hyperref}
\usepackage{fancyhdr}
\usepackage{titlesec}
\usepackage{xcolor}
\usepackage{setspace}

\definecolor{azulprincipal}{HTML}{2563EB}

% Interlineado 1.4
\setstretch{1.4}

% Espacio entre párrafos
\setlength{\parskip}{8pt}

\pagestyle{fancy}
\fancyhf{}
\renewcommand{\headrulewidth}{0pt}
\fancyfoot[C]{\thepage}

% Más espacio antes y después de secciones
\titleformat{\section}{\Large\bfseries\color{azulprincipal}}{\thesection.}{0.5em}{}
\titlespacing*{\section}{0pt}{28pt}{14pt}
\titleformat{\subsection}{\large\bfseries}{\thesubsection)}{0.5em}{}
\titlespacing*{\subsection}{0pt}{20pt}{10pt}

\begin{document}

% ─── INTRODUCCIÓN ───
\section*{Introducción}
\addcontentsline{toc}{section}{Introducción}

En este trabajo buscamos aplicar lo que vimos en Probabilidad y Estadística a un caso concreto: nuestro proyecto de Ingeniería de Software II. El sistema se llama BookDesk, y es una plataforma web para gestionar reservas de salas y escritorios en un espacio de coworking. Fue planteado por el profesor de la materia como trabajo práctico, y consiste en un sistema centralizado donde cada miembro puede ver la disponibilidad de los espacios en tiempo real y reservar desde el navegador.

Para el análisis estadístico nos enfocamos en la duración de las reservas como variable principal, porque nos pareció que es un dato útil para entender cómo se usa el espacio. También analizamos si hay relación entre la cantidad de usuarios conectados en un día y cuántas reservas se generan.

Trabajamos con datos simulados (50 observaciones para lo descriptivo y 30 para regresión) porque el sistema todavía no tiene suficiente historial real. Los datos los armamos tratando de que reflejen patrones de uso creíbles: reuniones cortas, jornadas parciales, jornadas completas, etc. Los resultados del análisis después los integramos en un módulo visual dentro de BookDesk para que el administrador pueda consultarlos.

% ═══════════════════════════════════════════════════════════════════════
\section{Recolección y organización de datos}
% ═══════════════════════════════════════════════════════════════════════

\subsection{Variable identificada}

\begin{itemize}[itemsep=4pt]
  \item Variable: Duración de las reservas (en horas).
  \item Tipo: Cuantitativa continua.
  \item Población: Todas las reservas que se hagan en BookDesk.
  \item Muestra: 50 observaciones simuladas.
\end{itemize}

Elegimos esta variable porque tiene impacto directo en cómo se administra el coworking. Si sabemos cuánto duran las reservas en promedio, podemos configurar mejor los slots de tiempo y anticipar cuándo van a quedar libres los espacios.

\subsection{Obtención de la muestra}

Generamos una muestra de $n = 50$ reservas simulando los distintos tipos de uso que esperamos en el sistema:

\begin{itemize}[itemsep=4pt]
  \item Reuniones cortas en salas (alrededor de 2 horas).
  \item Medias jornadas en escritorios (unas 4 horas).
  \item Jornadas completas en escritorios (entre 6 y 7 horas).
  \item Reservas rápidas de salas para llamadas o entrevistas (1 a 1,5 horas).
\end{itemize}

Los datos se mantienen fijos para que los resultados sean siempre reproducibles.

\subsection{Organización de datos}

Organizamos los datos con un histograma que muestra cómo se distribuyen las duraciones y un gráfico de torta para ver la proporción entre salas de reuniones y escritorios.

Lo que se ve en el histograma es que la mayoría de las reservas caen en el rango de 1 a 3 horas, y después hay un segundo grupo entre 5 y 7 horas que corresponde a la gente que reserva escritorios para todo el día. En cuanto al tipo de recurso, quedó bastante parejo: 25 reservas de salas y 25 de escritorios.

% ═══════════════════════════════════════════════════════════════════════
\section{Estadística descriptiva}
% ═══════════════════════════════════════════════════════════════════════

\subsection{Distribución de frecuencias}

Lo primero que hicimos fue armar la tabla de distribución de frecuencias. Para determinar el número de clases usamos la regla de Sturges:

\vspace{4pt}
$$k = 1 + 3{,}322 \cdot \log_{10}(n) = 1 + 3{,}322 \cdot \log_{10}(50) = 1 + 3{,}322 \cdot 1{,}699 = 6{,}644$$
\vspace{4pt}

Redondeando, nos quedan $k = 7$ clases.

Después calculamos la amplitud de cada clase. El valor mínimo de la muestra es $x_{min} = 0{,}50$ h y el máximo es $x_{max} = 7{,}91$ h, así que:

\vspace{4pt}
$$\text{Rango} = x_{max} - x_{min} = 7{,}91 - 0{,}50 = 7{,}41 \text{ h}$$

$$\text{Amplitud} = \frac{\text{Rango}}{k} = \frac{7{,}41}{7} = 1{,}059 \approx 1{,}1 \text{ h}$$
\vspace{4pt}

Arrancamos el primer intervalo en $0{,}5$ (el valor mínimo) y con amplitud $1{,}1$ nos quedan estos intervalos:

\vspace{6pt}
\begin{table}[H]
\centering
\caption{Distribución de frecuencias — Duración de reservas (horas)}
\vspace{4pt}
\begin{tabular}{lcccccc}
\toprule
\textbf{Clase} & \textbf{$x_i$} & \textbf{$f_i$} & \textbf{$F_i$} & \textbf{$h_i$} & \textbf{$H_i$} & \textbf{$h_i\%$} \\
\midrule
{[0,5 – 1,6)} & 1,05 & 7 & 7 & 0,1400 & 0,1400 & 14,00\% \\[3pt]
{[1,6 – 2,7)} & 2,15 & 17 & 24 & 0,3400 & 0,4800 & 34,00\% \\[3pt]
{[2,7 – 3,8)} & 3,25 & 6 & 30 & 0,1200 & 0,6000 & 12,00\% \\[3pt]
{[3,8 – 4,9)} & 4,35 & 7 & 37 & 0,1400 & 0,7400 & 14,00\% \\[3pt]
{[4,9 – 6,0)} & 5,45 & 5 & 42 & 0,1000 & 0,8400 & 10,00\% \\[3pt]
{[6,0 – 7,1)} & 6,55 & 5 & 47 & 0,1000 & 0,9400 & 10,00\% \\[3pt]
{[7,1 – 8,2)} & 7,65 & 3 & 50 & 0,0600 & 1,0000 & 6,00\% \\
\bottomrule
\end{tabular}
\end{table}
\vspace{4pt}

Para armar la tabla: las marcas de clase ($x_i$) son el punto medio de cada intervalo (por ejemplo, $(0{,}5 + 1{,}6)/2 = 1{,}05$). Las frecuencias relativas se calcularon como $h_i = f_i / n$ (por ejemplo, $17/50 = 0{,}34$), y las acumuladas se van sumando.

Se nota que el segundo intervalo (1,6 a 2,7 horas) concentra la mayor cantidad de reservas, con un 34\%. Tiene sentido porque muchas reuniones y medias jornadas duran entre 1,5 y 2,5 horas.

\subsection{Medidas de tendencia central}

Media aritmética: sumamos todas las duraciones y dividimos por $n$:

\vspace{4pt}
$$\bar{x} = \frac{\sum_{i=1}^{n} x_i}{n} = \frac{174{,}90}{50} = 3{,}4980 \text{ h}$$
\vspace{6pt}

Mediana: ordenamos las 50 observaciones de menor a mayor. Como $n = 50$ es par, la mediana es el promedio de las observaciones 25 y 26:

\vspace{4pt}
$$Me = \frac{x_{(25)} + x_{(26)}}{2} = 3{,}05 \text{ h}$$
\vspace{6pt}

Moda: buscamos el valor que más se repite. Redondeando a un decimal, la duración más frecuente es:

\vspace{4pt}
$$Mo = 2{,}0 \text{ h}$$

\subsection{Medidas de dispersión}

Varianza muestral: calculamos la suma de los desvíos al cuadrado y dividimos por $n - 1$:

\vspace{4pt}
$$s^2 = \frac{\sum_{i=1}^{n}(x_i - \bar{x})^2}{n - 1} = \frac{186{,}9654}{49} = 3{,}8156$$
\vspace{6pt}

Desvío estándar:

\vspace{4pt}
$$s = \sqrt{s^2} = \sqrt{3{,}8156} = 1{,}9534 \text{ h}$$
\vspace{6pt}

Coeficiente de variación:

\vspace{4pt}
$$CV = \frac{s}{\bar{x}} \times 100 = \frac{1{,}9534}{3{,}4980} \times 100 = 55{,}84\%$$

\subsection{Medidas de posición}

\begin{itemize}[itemsep=6pt]
  \item Primer cuartil: $Q_1 = 1{,}96$ h (el 25\% de las reservas dura menos que esto)
  \item Tercer cuartil: $Q_3 = 5{,}17$ h (el 75\% dura menos que esto)
  \item Rango intercuartílico: $RIQ = Q_3 - Q_1 = 5{,}17 - 1{,}96 = 3{,}21$ h
\end{itemize}

\subsection{Resumen de medidas}

\begin{table}[H]
\centering
\caption{Resumen de medidas estadísticas}
\vspace{4pt}
\begin{tabular}{lll}
\toprule
\textbf{Medida} & \textbf{Valor} & \textbf{Qué nos dice} \\
\midrule
Media ($\bar{x}$) & 3,50 h & El promedio de duración \\[4pt]
Mediana ($Me$) & 3,05 h & La mitad de las reservas dura menos que esto \\[4pt]
Moda ($Mo$) & 2,0 h & La duración que más se repite \\[4pt]
Desvío estándar ($s$) & 1,95 h & Cuánto varían respecto al promedio \\[4pt]
Coef. de variación ($CV$) & 55,84\% & Variabilidad alta \\[4pt]
$Q_1$ & 1,96 h & El 25\% más corto dura menos que esto \\[4pt]
$Q_3$ & 5,17 h & Solo el 25\% dura más que esto \\[4pt]
Rango & 7,41 h & Diferencia entre la más larga y la más corta \\
\bottomrule
\end{tabular}
\end{table}

\subsection{Análisis del comportamiento}

Lo primero que notamos es que la media (3,50h) es más alta que la mediana (3,05h). Esto indica que la distribución tiene asimetría positiva, o sea que hay algunas reservas largas que ``tiran'' el promedio para arriba. Tiene lógica: la mayoría de las reservas son cortas (reuniones, llamadas), pero las jornadas completas de 6-7 horas inflan la media.

El coeficiente de variación dio 55,84\%, que es bastante alto. Pero esto no es necesariamente malo para nuestro caso: es esperable que haya mucha variación porque el coworking se usa para cosas muy distintas. No es lo mismo una llamada de una hora que alguien que se sienta a trabajar todo el día.

Si miramos los cuartiles, el 50\% central de las reservas dura entre 1,96 y 5,17 horas. Esto nos da una idea de que el sistema debería manejar bien tanto reservas cortas como las de media jornada, que es donde cae la mayor parte del uso.

% ═══════════════════════════════════════════════════════════════════════
\section{Regresión y correlación}
% ═══════════════════════════════════════════════════════════════════════

\subsection{Variables seleccionadas}

Acá quisimos ver si hay relación entre cuántos usuarios se conectan al sistema en un día y cuántas reservas se generan ese mismo día:

\begin{itemize}[itemsep=4pt]
  \item $X$: Usuarios activos por día.
  \item $Y$: Reservas diarias.
\end{itemize}

Usamos datos de $n = 30$ días. La pregunta era bastante intuitiva: ¿si hay más gente conectada, se hacen más reservas? Parece obvio, pero queríamos cuantificarlo.

\subsection{Diagrama de dispersión}

Al graficar los puntos se ve una tendencia clara: a más usuarios, más reservas. Los puntos no están dispersos al azar sino que siguen una dirección ascendente bastante definida, así que decidimos ajustar una recta de regresión.

\subsection{Cálculos auxiliares}

Primero calculamos las sumatorias que vamos a necesitar para la regresión y la correlación:

\vspace{6pt}
\begin{table}[H]
\centering
\caption{Sumatorias para regresión y correlación ($n = 30$)}
\vspace{4pt}
\begin{tabular}{lr}
\toprule
\textbf{Sumatoria} & \textbf{Valor} \\
\midrule
$\sum x_i$ & 523 \\[4pt]
$\sum y_i$ & 400 \\[4pt]
$\sum x_i \cdot y_i$ & 8.275 \\[4pt]
$\sum x_i^2$ & 10.995 \\[4pt]
$\sum y_i^2$ & 6.302 \\
\bottomrule
\end{tabular}
\end{table}
\vspace{4pt}

Con eso sacamos las medias:

\vspace{4pt}
$$\bar{x} = \frac{\sum x_i}{n} = \frac{523}{30} = 17{,}4333 \qquad \bar{y} = \frac{\sum y_i}{n} = \frac{400}{30} = 13{,}3333$$

\subsection{Ecuación de regresión lineal}

Calculamos la pendiente $m$ y la ordenada al origen $b$ por mínimos cuadrados:

\vspace{4pt}
$$m = \frac{\sum_{i=1}^{n} x_i y_i - n \cdot \bar{x} \cdot \bar{y}}{\sum_{i=1}^{n} x_i^2 - n \cdot \bar{x}^2}$$
\vspace{4pt}

El numerador (covarianza):

$$\sum x_i y_i - n \cdot \bar{x} \cdot \bar{y} = 8275 - 30 \times 17{,}4333 \times 13{,}3333 = 1301{,}6667$$

El denominador (varianza de $x$):

$$\sum x_i^2 - n \cdot \bar{x}^2 = 10995 - 30 \times 17{,}4333^2 = 1877{,}3667$$

Entonces:

$$m = \frac{1301{,}6667}{1877{,}3667} = 0{,}6934$$

\vspace{6pt}
Y la ordenada al origen:

$$b = \bar{y} - m \cdot \bar{x} = 13{,}3333 - 0{,}6934 \times 17{,}4333 = 1{,}2459$$
\vspace{4pt}

La ecuación de la recta queda:

\vspace{4pt}
$$\boxed{\hat{y} = 0{,}6934 \cdot x + 1{,}2459}$$
\vspace{4pt}

Esto quiere decir que por cada usuario activo más que se conecta en un día, se esperan aproximadamente 0,69 reservas adicionales. El término independiente (1,25) representaría una especie de ``piso'' de reservas que se hacen incluso con muy pocos usuarios, probablemente los que ya tienen reservas programadas de antes.

\subsection{Coeficiente de correlación}

Calculamos el coeficiente de Pearson:

\vspace{4pt}
$$r = \frac{\sum x_i y_i - n \cdot \bar{x} \cdot \bar{y}}{\sqrt{\left(\sum x_i^2 - n \cdot \bar{x}^2\right) \cdot \left(\sum y_i^2 - n \cdot \bar{y}^2\right)}} = \frac{1301{,}6667}{\sqrt{1877{,}3667 \times 968{,}6933}} = \frac{1301{,}6667}{1348{,}5646} = 0{,}9652$$
\vspace{4pt}

Y el coeficiente de determinación:

\vspace{4pt}
$$r^2 = (0{,}9652)^2 = 0{,}9317$$
\vspace{4pt}

El $r$ de 0,97 es alto, lo que confirma que la relación entre usuarios activos y reservas es fuerte y positiva. Y el $r^2$ nos dice que el 93\% de la variación en las reservas se puede explicar por la cantidad de usuarios. Es un modelo que funciona bastante bien.

En la práctica, esto podría servir para que el administrador del coworking anticipe días de alta demanda. Si ve que hay muchos usuarios conectados temprano, ya sabe que va a haber muchas reservas y puede prepararse.

% ═══════════════════════════════════════════════════════════════════════
\section{Estadística inferencial}
% ═══════════════════════════════════════════════════════════════════════

\subsection{Intervalo de confianza}

Queríamos estimar cuál sería la duración media ``real'' de las reservas si tuviéramos datos de toda la población (no solo nuestra muestra de 50). Para eso armamos un intervalo de confianza al 95\%.

Datos:
\begin{itemize}[itemsep=4pt]
  \item $n = 50$, \quad $\bar{x} = 3{,}4980$ h, \quad $s = 1{,}9534$ h
  \item Nivel de confianza: $1 - \alpha = 0{,}95$ \quad $\Rightarrow$ \quad $\alpha = 0{,}05$
  \item Como no conocemos $\sigma$ poblacional, usamos la distribución $t$-Student
\end{itemize}

\vspace{6pt}
\textbf{Paso 1:} Grados de libertad:
$$gl = n - 1 = 50 - 1 = 49$$

\vspace{6pt}
\textbf{Paso 2:} Buscamos el valor crítico $t_{\alpha/2;\,gl}$ en la tabla de $t$-Student:
$$t_{0{,}025;\,49} = 2{,}0096$$

\vspace{6pt}
\textbf{Paso 3:} Calculamos el error estándar:
$$\hat{\sigma}_{\bar{x}} = \frac{s}{\sqrt{n}} = \frac{1{,}9534}{\sqrt{50}} = \frac{1{,}9534}{7{,}0711} = 0{,}2762$$

\vspace{6pt}
\textbf{Paso 4:} Calculamos el margen de error:
$$E = t_{\alpha/2} \cdot \hat{\sigma}_{\bar{x}} = 2{,}0096 \times 0{,}2762 = 0{,}5551$$

\vspace{6pt}
\textbf{Paso 5:} Armamos el intervalo:
$$IC = \bar{x} \pm E = 3{,}4980 \pm 0{,}5551$$

$$\boxed{IC_{95\%}(\mu) = (2{,}94 \text{ ; } 4{,}05) \text{ horas}}$$
\vspace{4pt}

Podemos decir con un 95\% de confianza que la duración promedio real de las reservas está entre 2 horas 56 minutos y 4 horas 3 minutos, aproximadamente. Para el administrador esto es útil porque le da un rango confiable para planificar la rotación de espacios.

\subsection{Prueba de hipótesis}

También nos preguntamos: ¿la duración promedio es distinta de 3 horas? Tres horas es un valor de referencia razonable (muchos coworkings configuran sus slots en bloques de 3h), así que planteamos:

\vspace{6pt}
\textbf{Paso 1 — Hipótesis:}
\begin{align*}
H_0&: \mu = 3 \text{ horas (la duración media es 3 horas)} \\
H_1&: \mu \neq 3 \text{ horas (la duración media es diferente de 3 horas)}
\end{align*}

Es una prueba bilateral (de dos colas) con $\alpha = 0{,}05$.

\vspace{6pt}
\textbf{Paso 2 — Estadístico de prueba:}

Usamos el estadístico $t$ porque no conocemos $\sigma$:

\vspace{4pt}
$$t_{calc} = \frac{\bar{x} - \mu_0}{s / \sqrt{n}} = \frac{3{,}4980 - 3{,}0}{1{,}9534 / \sqrt{50}} = \frac{0{,}4980}{0{,}2762} = 1{,}8027$$

\vspace{6pt}
\textbf{Paso 3 — Valor crítico y regla de decisión:}

Con $gl = 49$ y $\alpha/2 = 0{,}025$:
$$t_{crit} = \pm 2{,}0096$$

La regla es: si $|t_{calc}| > t_{crit}$, rechazamos $H_0$.

\vspace{6pt}
\textbf{Paso 4 — Decisión:}

Comparamos:
$$|t_{calc}| = 1{,}80 \quad < \quad t_{crit} = 2{,}01$$

El estadístico cae en la zona de no rechazo. También calculamos el $p$-valor:
$$p = 0{,}0777 \quad > \quad \alpha = 0{,}05$$

Como $p > \alpha$, no rechazamos $H_0$.

\vspace{6pt}
\textbf{Paso 5 — Conclusión:}

No tenemos evidencia estadística suficiente para decir que el promedio de duración sea distinto de 3 horas. El valor observado (3,50h) está por encima, pero la diferencia no es significativa al 5\%.

\subsection{¿Qué significan estos resultados juntos?}

Los dos resultados se complementan bien: el intervalo de confianza nos dice que la media real está entre 2,94 y 4,05 horas (y el valor 3 cae dentro de ese rango), y la prueba de hipótesis confirma que no hay diferencia significativa con respecto a 3 horas. En términos prácticos, los bloques de tiempo de 3 horas que usa el sistema están bien calibrados para el uso real.

% ═══════════════════════════════════════════════════════════════════════
\section{Integración con Ingeniería de Software II}
% ═══════════════════════════════════════════════════════════════════════

A partir de todo este análisis pudimos sacar varias conclusiones que impactan directamente en el desarrollo de BookDesk:

\begin{enumerate}[itemsep=10pt]
  \item \textbf{Slots de tiempo predeterminados:} Como las duraciones varían mucho ($CV = 55{,}8\%$), tiene sentido ofrecer opciones rápidas de 1h, 2h, 4h y jornada completa en vez de que el usuario tenga que elegir la hora de inicio y fin manualmente. Eso haría más ágil la reserva.

  \item \textbf{Anticipar la demanda:} El modelo de regresión con $r^2 = 0{,}93$ funciona bien para predecir reservas a partir de usuarios activos. Una idea sería agregar un indicador en el panel del administrador que muestre algo como ``hoy se esperan aproximadamente X reservas'' basándose en los usuarios que ya se conectaron.

  \item \textbf{Validar decisiones de diseño:} La prueba de hipótesis nos sirvió para confirmar algo que habíamos asumido al diseñar el sistema: que 3 horas es un buen valor por defecto para los bloques. Ahora tenemos un respaldo estadístico en vez de una suposición.

  \item \textbf{Detectar cambios en el uso:} El intervalo de confianza (2,94h – 4,05h) puede funcionar como una referencia. Si en el futuro el promedio de duración se sale de ese rango, puede ser señal de que algo cambió en cómo la gente usa el espacio y habría que revisar la configuración.

  \item \textbf{Módulo de estadística en el sistema:} Los resultados del análisis los cargamos en una sección del panel de administración de BookDesk (ruta \texttt{/admin/estadistica}), con gráficos y explicaciones simplificadas para que el administrador pueda consultarlos sin necesidad de saber estadística.
\end{enumerate}

% ─── CONCLUSIÓN ───
\section*{Conclusión}
\addcontentsline{toc}{section}{Conclusión}

Con este trabajo pudimos ver cómo las herramientas de estadística que aprendimos en la materia se pueden aplicar a un proyecto real (o al menos, a uno con datos simulados realistas). Logramos caracterizar el uso del coworking a través de las duraciones de reserva, encontrar una relación fuerte entre usuarios y demanda, y validar que la configuración actual del sistema tiene sentido estadístico.

Lo que más nos sirvió fue darnos cuenta de que los datos respaldan decisiones que habíamos tomado de forma intuitiva al diseñar BookDesk. Por ejemplo, los slots de 3 horas no los elegimos con ningún criterio formal, pero la prueba de hipótesis muestra que efectivamente es un buen valor.

También estuvo bueno poder volcar los resultados en el sistema como una sección del panel administrativo. La idea es que un administrador real pueda entrar y ver las métricas sin necesidad de saber estadística, con explicaciones claras de qué significa cada cosa.

\vspace{1cm}
\noindent\textbf{Herramientas utilizadas:}
\begin{itemize}[itemsep=4pt]
  \item Calculadora científica para los cálculos.
  \item React 18 con Recharts (librería de gráficos) para los histogramas, diagramas de dispersión y demás visualizaciones en el sistema.
  \item Supabase (PostgreSQL) como base de datos del sistema.
  \item Vercel para el despliegue de la aplicación.
\end{itemize}

\end{document}
